\documentclass[11pt]{article}
\usepackage[margin=0.85in]{geometry}
\usepackage{amsmath,amssymb,booktabs}
\usepackage[T1]{fontenc}
\usepackage{lmodern}
\usepackage[hidelinks]{hyperref}
\setlength{\parindent}{0pt}
\setlength{\parskip}{5pt}
\newcommand{\V}{\mathcal V}
\newcommand{\R}{\mathcal R}
\newcommand{\G}{\mathcal G}
\newcommand{\Q}{\mathcal Q}
\newcommand{\PP}{\mathcal P}
\begin{document}
\begin{center}
{\large\bfseries PS1, Part M.1: Discrete-State Formulation and Euler Equation}\\[3pt]
{\small ECONS 509 \quad Worked derivation for Question 2}
\end{center}

\textbf{1. Household problem and fixed prices.}
Let $k$ denote current assets, $k'$ next-period assets, and $s$ the labor-efficiency state.
The household solves
\begin{equation}
v(k,s)=\max_{k'\in\mathcal K,\ k'\geq\underline k,\ c>0}
\left\{u(c)+\beta\sum_{s'=1}^2\PP_{ss'}v(k',s')\right\},
\qquad c=(1+r^0)k+w^0\bar\epsilon_s\bar l-k',
\end{equation}
where $u(c)=c^{1-\sigma}/(1-\sigma)$ and $\underline k=0$.
Use $\beta=0.96$, $\sigma=1.5$, $\bar l=1$, $\bar\epsilon=(0.8,1.2)$,
and $\PP_{ss'}=0.5$ for every pair of shocks. Rows of $\PP$ condition on the current shock.
For $Y=zK^\alpha L^{1-\alpha}$, take $z=1$, $\alpha=0.36$, $\delta=0.1$, and $K^0=5$.
The shock probabilities are $(0.5,0.5)$, so $L^*=\sum_s\pi_s^{shock}\bar\epsilon_s\bar l=1$ and
\[
w^0=(1-\alpha)z(K^0/L^*)^\alpha\simeq1.1423762763,
\qquad r^0=\alpha z(K^0/L^*)^{\alpha-1}-\delta\simeq0.0285173311.
\]
These prices stay fixed; mean household assets need not equal $K^0$.

\textbf{2. Vector form.}
Let $\mathcal K=\{k_1,\ldots,k_N\}$, $S=2$, $M=NS$, and
$m(n,s)=(s-1)N+n$. Stack the columns of $\V$ in $x=\V(:)$:
all asset nodes in shock 1 precede those in shock 2.
\begin{center}
\begin{tabular}{lll}
\toprule
Object & Definition & Dimension\\
\midrule
$\V$ & $\V_{n,s}=v(k_n,s)$ & $N\times S$\\
$x=\V(:)$ & Column-stacked values & $M\times1$\\
$\R$ & Current-state and savings-choice utility & $M\times N$\\
$\G$ & Maximizing savings indices & $N\times S$\\
$\Q_G$ & Joint transition under policy $G$ & $M\times M$\\
$\R_G^*$, $\pi$ & Chosen utility, joint invariant distribution & $M\times1$ each\\
\bottomrule
\end{tabular}
\end{center}
Define $c_{n,s,n'}=(1+r^0)k_n+w^0\bar\epsilon_s\bar l-k_{n'}$ and
\[
\R_{m(n,s),n'}=
\begin{cases}
u(c_{n,s,n'}),& c_{n,s,n'}>0,\\
-\infty,&c_{n,s,n'}\leq0.
\end{cases}
\]
Since $(\PP\V^T)_{s,n'}=\sum_{s'}\PP_{ss'}\V_{n',s'}$, the stacked Bellman equation is
\begin{equation}
\boxed{\V(:)=\mathcal T(x)=\max_{\mathrm{columns}}
\left[\R+\beta(\PP\V^T)\otimes\mathbf1_N\right].}
\end{equation}
The maximization is over savings choices. Its argmax gives $\G(:)$, with the smallest
index selected at an exact tie. Set $g(k_n,s)=k_{\G(n,s)}$. Then
\begin{equation}
\Q_G[m(n,s),m(n',s')]=\PP_{ss'}\mathbf1\{n'=\G(n,s)\},
\qquad \R_G^*[m(n,s)]=\R[m(n,s),\G(n,s)].
\end{equation}
Each row of $\Q_G$ sums to one. For a fixed policy,
$x_G=\R_G^*+\beta\Q_Gx_G$; the column distribution satisfies
$\Q_G^T\pi=\pi$, $\mathbf1^T\pi=1$, and $\pi\geq0$.

\newpage
\textbf{3. Continuous-choice Euler equation.}
Now allow continuous $k'\geq\underline k$. With multiplier $\mu\geq0$ on
$k'-\underline k\geq0$, write
\[
\mathcal L=u\big((1+r^0)k+w^0\bar\epsilon_s\bar l-k'\big)
+\beta\sum_{s'}\PP_{ss'}v(k',s')+\mu(k'-\underline k).
\]
At an optimum, the first-order and complementary-slackness conditions are
\begin{align}
-u'(c)+\beta\sum_{s'}\PP_{ss'}v_k(k',s')+\mu&=0,\\
\mu\geq0,\qquad k'-\underline k\geq0,\qquad\mu(k'-\underline k)&=0.
\end{align}
The envelope condition, holding prices fixed, is
\begin{equation}
v_k(k,s)=(1+r^0)u'(c(k,s)).
\end{equation}
Evaluating the envelope condition at next period's state and using $u'(c)=c^{-\sigma}$ gives
\begin{equation}
\boxed{c(k,s)^{-\sigma}
=\beta(1+r^0)\sum_{s'}\PP_{ss'}c(k',s')^{-\sigma}+\mu.}
\end{equation}
Thus current marginal utility is at least discounted expected next-period marginal utility
times the gross return. If $k'>\underline k$, then $\mu=0$ and the Euler equation holds
with equality. At the borrowing limit, it may be a strict inequality. This condition
assumes no active artificial upper asset bound; choices at that grid boundary are flagged
as a possible truncation issue.

\textbf{4. Euler residual for the discrete policy.}
Using the grid policy, compute consumption by exact index lookup:
\begin{align}
c_{n,s}&=(1+r^0)k_n+w^0\bar\epsilon_s\bar l-k_{\G(n,s)},\\
c^{next}_{n,s,s'}&=(1+r^0)k_{\G(n,s)}+w^0\bar\epsilon_{s'}\bar l
-k_{\G(\G(n,s),s')}.
\end{align}
Define the Euler-implied consumption and the absolute, unit-free consumption error as
\begin{equation}
c^{EE}_{n,s}=\left[\beta(1+r^0)\sum_{s'}\PP_{ss'}
(c^{next}_{n,s,s'})^{-\sigma}\right]^{-1/\sigma},
\qquad \boxed{E_{n,s}=\left|1-\frac{c^{EE}_{n,s}}{c_{n,s}}\right|.}
\end{equation}
For part (h), use all states in $A=\{(n,s):k_{\G(n,s)}>0\}$, flag upper-bound choices,
and plot $E_{n,s}$ against $k_n$ for each shock. Report the maximum and arithmetic mean
over $A$. With conditional weights $q_{n,s}=\pi_{n,s}/\sum_A\pi_{n,s}$,
report $\sum_Aq_{n,s}E_{n,s}$ and $\max_A(q_{n,s}E_{n,s})$, labeling the latter
the maximum weighted contribution. Also report the maximum over slack states with
$\pi_{n,s}>10^{-12}$. Empty sets or zero weight denominators give unavailable summaries.
Discrete choices need not satisfy the continuous Euler equality exactly.

\textbf{5. Link to the M.2 kernel.}
The kernel uses $N=100$ uniform points on $[0,20]$, starts from $\V_0=0$, and stops when
\[
\max_m\frac{|\mathcal T(x)_m-x_m|}{|x_m|+1}\leq10^{-10}.
\]
It returns the current values and their greedy policy before another update; the iteration
count includes the stopping check. Power iteration starts at $\pi_0=\mathbf1/M$ and returns
the updated distribution once $\|\pi_{j+1}-\pi_j\|_\infty\leq10^{-12}$.
In Python, column stacking uses \texttt{order="F"}, and saved policy indices are zero-based.
The script saves \texttt{V}, \texttt{G}, \texttt{pi}, and \texttt{iterations} in
\texttt{kernel\_output.npz}, matching the specification's interface.
\end{document}
